\section{Exact output}\label{sec:exact}

A certified gap below $2^{-q}$ is not an exact solution. For rational
quadratics, a global minimizer and the optimal value have bounded rational
height, so a sufficiently small certified gap can be converted into an exact
optimizer that is checked against the original data. This section proves a
height bound for minimizers and for the optimal value
(Section~\ref{sec:exact-height}), gives an acceptance rule that needs neither
uniqueness nor growth, and shows that an exact optimizer can be recovered from
any feasible point close to the \emph{set} of minimizers
(Section~\ref{sec:exact-accept}). Under point growth this yields exact output
in $f_1(p,\kappa)\poly(I)$ bit operations (Theorem~\ref{thm:exact}).
Section~\ref{sec:localized} gives an acceptance test that uses the filtering
history, and Section~\ref{sec:polynomial} treats explicit polynomial factors.
Several minimizers are treated in Section~\ref{sec:optsets}.

\subsection{Scaled data}\label{sec:exact-data}

In this section, except in Section~\ref{sec:polynomial}, $F$ is a quadratic
with rational data,
\[
F(x)=\tfrac12x^{\T}Hx+b^{\T}x+c ,
\]
where $H$ is the symmetric Hessian, on a bounded mixed box
$X=\prod_iX_i$ with rational endpoints, after fixed coordinates have been
substituted and integer endpoints rounded inward; thus $\ell_i<u_i$ for every
$i$. The curvature bounds are $L_i=\max\{H_{ii},0\}$, so
$P=\{i:H_{ii}>0\}$. Let $\Delta$ be a positive common denominator of the
monomial coefficients of $F$, namely $H_{ii}/2$, $H_{ij}$ ($i<j$), $b_i$ and
$c$, and of all endpoints $\ell_i,u_i$. Put
\begin{equation}\label{eq:exact-constants}
\hat H=\Delta H,\qquad I_C^+=I_C\cap P,\qquad
R=\Delta\prod_{i\in I_C^+}\hat H_{ii},\qquad \Omega=\Delta R^2,\qquad
\tau=\frac1{4nR}.
\end{equation}
The matrix $\hat H$ is integral and symmetric, since
$\hat H_{ii}=2\Delta(H_{ii}/2)$ and $\hat H_{ij}=\Delta H_{ij}$; this choice of
$\Delta$ absorbs the factor two of the off-diagonal entries. Each
$\hat H_{ii}$ with $i\in I_C^+$ is a positive integer, and
$\hat H_{ii}=\Delta L_i\le\Delta L$, so
$\log_2R\le\log_2\Delta+\abs{I_C^+}\log_2(\Delta L)$ is polynomial in $I$.
Every positive $H_{ii}$ is an integer multiple of $2/\Delta$, so $L\ge2/\Delta$
if $L>0$. The numbers $R$, $\Omega$ and $\tau$ are computed from the original
data; refined boxes never enter them.

Two elementary facts are used repeatedly. If every coordinate of $x\in X$
lies in $\rho^{-1}\Z$ for a positive integer $\rho$, then
$\Delta\rho^2F(x)\in\Z$, because $\Delta$ clears every monomial coefficient. If
$a/W\ne a'/W'$ are rationals with positive denominators, then
$\abs{a/W-a'/W'}\ge1/(WW')$.

\begin{lemma}[Hadamard's inequality {\cite[Theorem~7.8.1]{HornJohnson2013}}]\label{lem:hadamard}
If $A$ is a real symmetric positive definite $m\times m$ matrix, then
$\det A\le\prod_{i=1}^mA_{ii}$.
\end{lemma}

\subsection{Stationary polytopes and rational height}\label{sec:exact-height}

The argument follows Vavasis \cite{Vavasis1990} (see also
\cite{DelPiaDeyMolinaro2017}); we need its explicit constants. For $x\in X$ let
$J_0(x)=\{i\in I_C:\ell_i<x_i<u_i\}$ be its interior coordinates.

\begin{definition}[Stationary polytope]\label{def:statpoly}
For $s\in\mathcal S$ let
\[
\mathcal P(s)=\bigl\{x\in\R^n:\ x_i=s_i\ (i\notin J_0(s)),\ \
\ell_i\le x_i\le u_i\ \text{and}\ \partial_iF(x)=0\ (i\in J_0(s))\bigr\}.
\]
\end{definition}

\begin{lemma}[Stationary polytopes]\label{lem:statpoly}
Let $s\in\mathcal S$ and $J_0=J_0(s)$.
\begin{enumerate}[label=(\alph*)]
\item $\nabla_{J_0}F(s)=0$ and $H_{J_0J_0}\succeq0$.
\item $\mathcal P(s)$ is a nonempty polytope contained in $X$, and
$\mathcal P(s)\subseteq\mathcal S$.
\item Let $v$ be a vertex of $\mathcal P(s)$ and $J_v=J_0(v)$. Then
$J_v\subseteq J_0\cap I_C^+$, $\hat H_{J_vJ_v}\succ0$, every coordinate of $v$ lies
in $\rho^{-1}\Z$ for $\rho=\Delta\det\hat H_{J_vJ_v}$, and $\Delta\le\rho\le R$
(with $\det\hat H_{J_vJ_v}=1$ if $J_v=\emptyset$).
\end{enumerate}
\end{lemma}

\begin{proof}
(a) For $w\in\R^n$ supported on $J_0$, the point $s+tw$ lies in $X$ for all
small $\abs t$, because the coordinates in $J_0$ are continuous and strictly
inside their bounds. Since
$F(s+tw)=\OPT+t\nabla F(s)^{\T}w+\frac{t^2}2w^{\T}Hw$ has a minimum at $t=0$,
we get $\nabla F(s)^{\T}w=0$ and $w^{\T}Hw\ge0$.

(b) By (a), $s\in\mathcal P(s)$. The set is defined by finitely many linear
equations and inequalities inside a bounded box, so it is a polytope. Its
points keep the integer coordinates of $s$ and satisfy all bounds, so
$\mathcal P(s)\subseteq X$. For $x\in\mathcal P(s)$, $d=x-s$ is supported on
$J_0$, and $H_{J_0J_0}d_{J_0}=\nabla_{J_0}F(x)-\nabla_{J_0}F(s)=0$. Hence
\[
F(x)-F(s)=\nabla F(s)^{\T}d+\tfrac12d^{\T}Hd
=\nabla_{J_0}F(s)^{\T}d_{J_0}+\tfrac12d_{J_0}^{\T}H_{J_0J_0}d_{J_0}=0 ,
\]
so $x\in\mathcal S$.

(c) A continuous coordinate outside $J_0$ is fixed at $s_i\in\{\ell_i,u_i\}$
in $\mathcal P(s)$, so $J_v\subseteq J_0$. Suppose $w\in\R^{J_v}$ is nonzero and
$H_{J_0J_v}w=0$. Extend $w$ by zero to $\tilde w\in\R^n$. Then
$v\pm t\tilde w\in\mathcal P(s)$ for small $t>0$: the equations are preserved,
the coordinates in $J_v$ have slack, and all other coordinates are unchanged.
This contradicts that $v$ is a vertex, so $H_{J_0J_v}$ has trivial kernel. The
matrix $H_{J_vJ_v}$ is a principal submatrix of $H_{J_0J_0}\succeq0$; if
$w\in\R^{J_v}\setminus\{0\}$ had $w^{\T}H_{J_vJ_v}w=\tilde w^{\T}H\tilde w=0$, then
$H_{J_0J_0}\tilde w_{J_0}=0$, that is $H_{J_0J_v}w=0$. Hence $H_{J_vJ_v}\succ0$ and
$\hat H_{J_vJ_v}=\Delta H_{J_vJ_v}\succ0$. In particular $H_{ii}>0$ for $i\in J_v$, so
$J_v\subseteq I_C^+$.

The coordinates of $v$ outside $J_v$ are integers or endpoints, so they lie in
$\Delta^{-1}\Z$. The rows $i\in J_v$ of $\nabla F(v)=0$, multiplied by
$\Delta^2$, read
\[
\hat H_{J_vJ_v}(\Delta v_{J_v})=-\Delta(\Delta b_{J_v})-\sum_{j\notin J_v}\hat H_{J_vj}(\Delta v_j),
\]
whose right-hand side is integral. Cramer's rule gives
$\Delta v_{J_v}\in(\det\hat H_{J_vJ_v})^{-1}\Z^{J_v}$, so every coordinate of $v$ lies in
$\rho^{-1}\Z$. Finally $\det\hat H_{J_vJ_v}$ is a positive integer, and by
Lemma~\ref{lem:hadamard} and $J_v\subseteq I_C^+$,
$\det\hat H_{J_vJ_v}\le\prod_{i\in J_v}\hat H_{ii}\le\prod_{i\in I_C^+}\hat H_{ii}$.
\end{proof}

\begin{corollary}[Rational height]\label{cor:height}
(a) Some global minimizer has all coordinates in $\rho^{-1}\Z$ for an integer
$\rho$ with $\Delta\mid\rho$ and $\rho\le R$; every isolated global minimizer
has this property. The endpoints $\ell_i,u_i$ then also lie in $\rho^{-1}\Z$.
(b) $\OPT=a/W$ in lowest terms with $1\le W\le \Omega$.
\end{corollary}

\begin{proof}
(a) Take any $s\in\mathcal S$. The nonempty polytope $\mathcal P(s)$ has a
vertex, which is a global minimizer by Lemma~\ref{lem:statpoly}(b) and has
the stated height by Lemma~\ref{lem:statpoly}(c). If $s$ is isolated in
$\mathcal S$, then the convex set $\mathcal P(s)\subseteq\mathcal S$ contains
no other point near $s$, hence equals $\{s\}$, and $s$ is its vertex. The
endpoints lie in $\Delta^{-1}\Z\subseteq\rho^{-1}\Z$.
(b) Evaluate $F$ at the minimizer of (a): $\Delta\rho^2\OPT\in\Z$, so
$W\le \Delta\rho^2\le \Omega$.
\end{proof}

\begin{remark}[Other admissible height constants]\label{rem:heights}
Any bound on the minors of $\hat H$ restricted to $I_C$ also works in
Lemma~\ref{lem:statpoly}(c) and Corollary~\ref{cor:height}, for example
$\abs{I_C}!\,\max\{1,\max_{ij}\abs{\hat H_{ij}}\}^{\abs{I_C}}$ by the Leibniz
expansion, or $\prod_{i\in I_C}\max\{1,\sum_{j\in I_C}\abs{\hat H_{ij}}\}$ by
Hadamard's inequality for rows. Both also bound \emph{nonprincipal} minors,
which arise if a vertex of $\mathcal P(s)$ is described by a nonprincipal set
of stationarity rows. Part~(c) shows that the principal system
$\hat H_{J_vJ_v}$ always suffices, and Lemma~\ref{lem:hadamard} then gives the
diagonal bound in~\eqref{eq:exact-constants}, which is never larger than
either alternative and depends only on the coordinate curvatures. All these
constants have polynomial bit length; only the threshold values below
change.
\end{remark}

\subsection{Acceptance and recovery}\label{sec:exact-accept}

\begin{proposition}[Exact acceptance]\label{prop:accept}
Let $F$ attain its minimum $\OPT$ over the feasible set $X$, and let $\OPT$ be
a rational number with denominator at most $\Omega'$. Let $\beta\le\OPT$ and
let $x\in X$ with $F(x)=a/W$ in lowest terms. If
$F(x)-\beta<1/(\Omega'W)$, then $x$ is a global minimizer. For the quadratics
of this section, $\Omega'=\Omega$ is admissible by
Corollary~\ref{cor:height}(b).
\end{proposition}

\begin{proof}
If $F(x)>\OPT$, then $F(x)-\OPT\ge1/(\Omega'W)$ by the second elementary
fact, because the denominator of $\OPT$ is at most $\Omega'$. With
$\beta\le\OPT$ this contradicts $F(x)-\beta<1/(\Omega'W)$.
\end{proof}

The rule uses the denominator of the candidate itself. It does not assume
that feasible values lie on a lattice, which is false when continuous
coordinates vary. The proof uses only the rationality of $F(x)$ and the
denominator bound for $\OPT$; Sections~\ref{sec:recourse}
and~\ref{sec:constraints} apply the proposition with other feasible sets and
other values of $\Omega'$.

\begin{algorithm}[REC: snapping recovery]\label{alg:rec}
The input is a point $y\in X$. Fix every integer coordinate at $y_i$. For
every continuous coordinate, fix it at $\ell_i$ if $y_i-\ell_i\le\tau$, else
at $u_i$ if $u_i-y_i\le\tau$, and otherwise call it free; let $J$ be the set
of free coordinates. Find by exact linear programming a point $\tilde x$ that
agrees with the fixed values and satisfies $\ell_J\le\tilde x_J\le u_J$ and
$\nabla_JF(\tilde x)=0$, or report failure.
\end{algorithm}

The data of this linear system are original coefficients, endpoints and
integer values of $y$, so it is solved in time polynomial in $I$
\cite[Theorem~6.4.12]{GrotschelLovaszSchrijver1988}. If $H_{JJ}$ is
nonsingular, Gaussian elimination suffices.

\begin{lemma}[Snapping recovery]\label{lem:snap}
If $y\in X$ and $\dist(y,\mathcal S)\le\tau/2$, then REC$(y)$ does not fail,
and every point it can return lies in $\mathcal S$.
\end{lemma}

\begin{proof}
Let $s\in\mathcal S$ be nearest to $y$; it exists since $\mathcal S$ is
compact. As $\norm{y-s}\le\tau/2<1$, the integer coordinates of $y$ and $s$
agree. The two endpoints of a continuous coordinate differ by at least
$1/\Delta\ge1/R=4n\tau>2\tau$, so no coordinate is within $\tau$ of both
endpoints. If $s_i=\ell_i$ then $y_i-\ell_i\le\tau/2$ and coordinate $i$ is
fixed at $\ell_i$; similarly for $s_i=u_i$. Hence every continuous coordinate
at a bound in $s$ is fixed at its value in $s$, and $J\subseteq J_0=J_0(s)$.
Let $J'=J_0\setminus J$ be the remaining fixed coordinates, each fixed at an
endpoint $e_i\in\{\ell_i,u_i\}$, and define on $\mathcal P(s)$ the linear
function
\[
\phi(x)=\sum_{i\in J'}\abs{x_i-e_i}
=\sum_{i\in J',\,e_i=\ell_i}(x_i-\ell_i)+\sum_{i\in J',\,e_i=u_i}(u_i-x_i)\ge0 .
\]
For $i\in J'$ we have $\abs{y_i-e_i}\le\tau$ and $\abs{s_i-y_i}\le\tau/2$, so
$\phi(s)\le\frac32\abs{J'}\tau\le\frac3{8R}<\frac1R$. The minimum of $\phi$
over the polytope $\mathcal P(s)$ is attained at a vertex $v$. By
Lemma~\ref{lem:statpoly}(c) the coordinates of $v$ and all endpoints lie in
$\rho^{-1}\Z$ for an integer $\rho\le R$, so $\phi(v)\in\rho^{-1}\Z$.
Since $0\le\phi(v)\le\phi(s)<1/R\le1/\rho$, we get $\phi(v)=0$.

Thus $s'=v\in\mathcal P(s)\subseteq\mathcal S$ satisfies $s'_i=e_i$ for
$i\in J'$, agrees with $s$ (hence with the fixed values of REC) outside $J_0$,
and has $\nabla_JF(s')=0$ because $J\subseteq J_0$. So the system solved by
REC is feasible. If $\tilde x$ is any solution, then $d=\tilde x-s'$ is
supported on $J$ and $H_{JJ}d_J=\nabla_JF(\tilde x)-\nabla_JF(s')=0$, so
$F(\tilde x)-F(s')=\nabla_JF(s')^{\T}d_J+\frac12d_J^{\T}H_{JJ}d_J=0$.
Since $\tilde x\in X$, it is a global minimizer. No nonsingularity of
$H_{JJ}$ is used.
\end{proof}

\begin{theorem}[From certified gaps to exact minimizers]\label{thm:transfer}
Let $\beta\le\OPT$ be certified and $y\in X$.
\begin{enumerate}[label=(\alph*)]
\item If $F$ has set growth~\eqref{eq:setgrowth} with constant $g_S$ and
\begin{equation}\label{eq:transfer-threshold}
F(y)-\beta\le\varepsilon_S:=\min\Bigl\{\frac{g_S\tau^2}{4},\
\frac1{2\Omega^2}\Bigr\}
=\min\Bigl\{\frac{g_S}{64n^2R^2},\ \frac1{2\Omega^2}\Bigr\},
\end{equation}
then REC$(y)$ returns a point $\tilde x$ that passes the test of
Proposition~\ref{prop:accept} with $\Omega'=\Omega$.
\item Without any growth assumption there is $\varepsilon_0>0$, depending on
the instance, such that the conclusion of (a) holds whenever
$F(y)-\beta\le\varepsilon_0$.
\end{enumerate}
\end{theorem}

\begin{proof}
(a) $g_S\dist(y,\mathcal S)^2\le F(y)-\OPT\le F(y)-\beta\le g_S\tau^2/4$, so
Lemma~\ref{lem:snap} applies and $\tilde x\in\mathcal S$. Then
$F(\tilde x)=\OPT=a/W$ with $W\le \Omega$, and
$F(\tilde x)-\beta\le F(y)-\beta\le1/(2\Omega^2)<1/(\Omega W)$.
(b) The set $Y=\{x\in X:\dist(x,\mathcal S)\ge\tau/2\}$ is compact and
$F>\OPT$ on it, so $\varepsilon_Y=\min_YF-\OPT>0$ (put $\varepsilon_Y=\infty$
if $Y=\emptyset$). With $\varepsilon_0=\min\{\varepsilon_Y/2,1/(2\Omega^2)\}$,
$F(y)-\beta\le\varepsilon_0$ gives $\dist(y,\mathcal S)<\tau/2$, and the
argument of (a) applies.
\end{proof}

\begin{remark}[Set growth is automatic]\label{rem:setgrowth}
Every rational mixed box QP has set growth with some constant $g_S>0$. On an
integer slice whose continuous box contains a minimizer, Luo and Sturm's error
bound for a quadratic on a polytope \cite[Theorem~3.3]{LuoSturm2000}, applied
to $F-\OPT$, gives
$F(x)-\OPT\ge c\,\dist(x,\mathcal S\cap\text{slice})^2\ge c\,\dist(x,\mathcal S)^2$;
on the finitely many other slices $F-\OPT$ is bounded below by a positive
constant while $\dist(x,\mathcal S)$ is bounded by the diameter of $\bar X$.
As for point growth, no useful lower bound on $g_S$ follows.
Theorem~\ref{thm:transfer} shows that exactness costs only
$\log_2(1/\varepsilon_S)\le\poly(I)+\max\{0,\log_2(1/g_S)\}$ bits of
certified accuracy; the cost of reaching that accuracy is a property of the
approximation algorithm.
\end{remark}

\begin{remark}[Relation to rational-witness results]\label{rem:np}
Corollary~\ref{cor:height} is the box case, with explicit constants, of a
classical fact: a quadratic program with rational data that attains its
minimum has a global minimizer of polynomial encoding length
\cite{Vavasis1990}. Vavasis used it to show that the decision version of
quadratic programming is in NP, and Del Pia, Dey and Molinaro extended NP
membership to mixed-integer quadratic programs with unbounded integer
variables \cite{DelPiaDeyMolinaro2017}; on a bounded mixed box the integer
coordinates have polynomial length, so that extension is not needed. These
results certify $\OPT\le t$. Deciding $\OPT\le t$ for box QP is NP-complete,
and NP-hard already at treewidth two \cite{DelPiaKhajavirad2026}, so
certificates of $\OPT>t$ of polynomial size, checkable in polynomial time, for
all instances would place an NP-complete problem in coNP. The certificates of
Theorem~\ref{thm:exact} below have size $f_1(p,\kappa)\poly(I)$: they exist
for bounded $(p,\kappa)$, and they are checked without trusting growth. The
explicit constants $R$ and $\Omega$ are what make the stopping test
computable.
\end{remark}

\subsection{Exact output under point growth}\label{sec:exact-growth}

\begin{lemma}[Uniqueness implies growth]\label{lem:unique-growth}
If $\mathcal S=\{x^*\}$, then $F$ has point growth~\eqref{eq:growth} with
some constant $g>0$.
\end{lemma}

For quadratic programs over polyhedra this is part of the second-order theory
of \cite{Contesse1980}; the proof below also covers integer coordinates.

\begin{proof}
Otherwise there are $x^k\in X\setminus\{x^*\}$ with
$(F(x^k)-\OPT)/\norm{x^k-x^*}^2\to0$. As $\bar X$ is bounded,
$F(x^k)\to\OPT$, and compactness of $X$ with uniqueness gives $x^k\to x^*$;
so the integer coordinates of $x^k$ equal those of $x^*$ for large $k$, and
since $x^k\ne x^*$, they differ in some continuous coordinate. Put
$t_k=\norm{x^k-x^*}$ and $u^k=(x^k-x^*)/t_k$; pass to a subsequence with
$u^k\to u$, $\norm u=1$, $u_{I_Z}=0$. Then
\[
\frac{F(x^k)-\OPT}{t_k^2}=\frac{\nabla F(x^*)^{\T}u^k}{t_k}
+\frac12(u^k)^{\T}Hu^k .
\]
Since $x^*$ minimizes $F$ over the convex continuous box of its integer slice
and $x^k$ lies in that box, the first term is nonnegative. The left side
tends to zero and the second term is bounded, so
$\nabla F(x^*)^{\T}u=0$ and $u^{\T}Hu\le0$. Each $u^k$ satisfies
$u^k_i\ge0$ when $x^*_i=\ell_i$ and $u^k_i\le0$ when $x^*_i=u_i$; so does
$u$, and hence $x^*+tu\in X$ for small $t>0$. Then
$F(x^*+tu)-\OPT=\frac{t^2}2u^{\T}Hu\le0$, contradicting uniqueness.
\end{proof}

As noted after Definition~\ref{def:growth}, the lemma gives no bound on the
condition number, which can be exponential in $I$.

\begin{algorithm}[EX: exact output]\label{alg:ex}
For $q=1,2,4,8,\dots$: run CT$(2^{-q})$ (Algorithm~\ref{alg:ct}) and obtain a
feasible $\hat x$ and a certified $\beta$; run REC$(\hat x)$; if it returns
$\tilde x$ with $F(\tilde x)-\beta<1/(\Omega W)$, where $W$ is the reduced
denominator of $F(\tilde x)$, stop and return $\tilde x$, $F(\tilde x)$ and
the path certificate of CT.
\end{algorithm}

If $L=0$, then $P=\emptyset$, CT is exact and returns a vertex $\hat x$ of the
box with $F(\hat x)=\beta$; REC fixes every coordinate of a vertex, so
$\tilde x=\hat x$ and the test passes at $q=1$.

\begin{theorem}[Exact rational output]\label{thm:exact}
For every rational mixed box QP with a supplied tree decomposition, EX
stops, and its output is a global minimizer and $\OPT$, with a certificate
whose validity uses neither growth nor uniqueness. If the instance has point
growth~\eqref{eq:growth} with condition number $\kappa$, EX uses at most
$f_1(p,\kappa)(I+1)^{C_1}$ bit operations, where $f_1$ is a computable
function and $C_1$ an absolute constant; neither $g$ nor $\kappa$ is an
input. This bound remains valid if REC skips every call whose matrix $H_{JJ}$
is singular, so under point growth no linear programming is needed.
\end{theorem}

\begin{proof}
\emph{Validity.} Each returned point passes Proposition~\ref{prop:accept}
with a certified $\beta$ (Theorem~\ref{thm:approx}), so it is a global
minimizer.

\emph{Termination.} CT stops on every instance
(Theorem~\ref{thm:approx}(a)). Once $2^{-q}\le\varepsilon_0$ of
Theorem~\ref{thm:transfer}(b), the test passes.

\emph{Bound under growth.} If $L=0$, EX stops after one call of CT. Otherwise
$L\ge2/\Delta$. Point growth is set growth with $\mathcal S=\{x^*\}$ and
$g_S=g$. By Theorem~\ref{thm:transfer}(a), EX stops at the latest at the first
$q$ with $2^{-q}\le\varepsilon_S$, so the largest $q$ used is less than $2q^*$
with $q^*=\lceil\log_2(1/\varepsilon_S)\rceil\ge1$. Since $g\ge L/\kappa$ and
$\log R$, $\log \Omega$, $\log n$ and $\log(1/L)$ are at most polynomial in
$I$, $q^*\le\poly(I)+\log_2\kappa$. There are $O(\log q^*)$ calls of CT, each
within $f(p,\kappa)(I+2q^*+1)^5$ bit operations by Theorem~\ref{thm:approx},
whose bound holds with $\kappa\ge\bar\kappa$, and
$(I+2q^*+1)^5\le(1+\log_2\kappa)^5\poly(I)$. REC, the evaluation of
$F(\tilde x)$ and the comparison cost $\poly(I+q^*)$. Absorbing all powers
of $1+\log_2\kappa$ into $f_1$ proves the bound.

Finally, let $J_0=J_0(x^*)$. By Lemma~\ref{lem:statpoly}(a),
$\nabla_{J_0}F(x^*)=0$, so a nonzero $w$ supported on $J_0$ with
$H_{J_0J_0}w_{J_0}=0$ would make $F$ constant on a short segment through
$x^*$; hence $H_{J_0J_0}\succ0$. If $\dist(\hat x,x^*)\le\tau/2$, then
$J\subseteq J_0$ as in the proof of Lemma~\ref{lem:snap}, so $H_{JJ}$ is
nonsingular. This holds at every call that the bound uses, so skipping
singular calls does not change it.
\end{proof}

\begin{remark}[Coordinatewise reconstruction]\label{rem:cf}
Under point growth one may instead round each coordinate of $\hat x$ to the
unique rational of denominator at most $R$ within $1/(4R^2)$ (continued
fractions, \cite[Theorem~5.1.9]{GrotschelLovaszSchrijver1988}); this succeeds
once $2^{-q}\le g/(32R^4)$, by Corollary~\ref{cor:height}(a) and the
$1/R^2$ separation of such rationals. Snapping recovery succeeds once
$2^{-q}\le g/(64n^2R^2)$ (Lemma~\ref{lem:snap}), which is the weaker
requirement when $R^2\ge2n^2$; it does not need uniqueness, and it returns a minimizer even
when $\mathcal S$ is a continuum. In both routes, $2^{-q}\le1/(2\Omega^2)$
makes the recovered minimizer pass the acceptance test.
\end{remark}

\subsection{Localized acceptance}\label{sec:localized}
The acceptance rule of Proposition~\ref{prop:accept} needs a certified gap
below $1/(\Omega W)$, as small as about $2^{-315}$ on the instances of
experiment E4 (Section~\ref{sec:comp-exact}), which have at most six
variables. The filtering
history permits a different test that accepts a candidate as soon as the retained box
is small compared with a first-order margin.

Consider stage $j$ of a filtering run, such as a trial of CT: a grid $G$ for
the current box $X^{(j)}$ with corrections $d_i(v)=L_iw_i(v)^2/8$ and the
grids $G_i=\{\ell_i,u_i\}$ for $i\notin P$, min-marginals $m_i$, a corrected
minimizer $y_j$, the incumbent value $U$, and the filtered box $X^{(j+1)}$.
The boxes are reached from $X$ by filtering steps whose thresholds are values
of feasible points and at least $U$, as in a trial, where the incumbent value
never increases. So by Proposition~\ref{prop:filter} every $x\in X$ with
$F(x)\le U$ lies in $X^{(j+1)}$, and $y_j\in X^{(j+1)}$. Filtering never
narrows a coordinate $i\notin P$, so $X^{(j+1)}_i=X^{(j)}_i=X_i$.

\begin{lemma}[Node exclusion]\label{lem:node}
Let $x\in X^{(j)}$ with $F(x)\le U$.
\begin{enumerate}[label=(\alph*)]
\item If $i\in I_Z$, then $x_i$ lies in the set $Z_i$ of nodes $v$ with
$m_i(v)\le U$ and of integers strictly inside grid intervals $[a,a']$ with
$a'-a\ge2$ and $\min\{m_i(a),m_i(a')\}\le U$.
\item If $i\notin P$, $G_i=\{a,a'\}$ and $m_i(a)>U$, then setting $x_i=a'$
gives a point of $X^{(j)}$ with value at most $F(x)$.
\end{enumerate}
\end{lemma}

\begin{proof}
(a) If $x_i$ is a node $v$, Proposition~\ref{prop:cellwise}(c) gives
$m_i(v)\le m_i(v)+d_i(v)\le F(x)\le U$. Otherwise $x_i$ is an integer strictly
inside a grid interval $[a,a']$, which then has $a'-a\ge2$, and
Proposition~\ref{prop:cellwise}(c) gives $\min\{m_i(a),m_i(a')\}\le F(x)\le U$.
(b) Since $L_i=0$, $F$ is concave along coordinate $i$, so $F(x)$ is at least
the smaller of its values at $x_i=a$ and $x_i=a'$. By
Proposition~\ref{prop:cellwise}(c) the first is at least $m_i(a)>U\ge F(x)$,
so the second is at most $F(x)$.
\end{proof}

\begin{proposition}[Localized exact acceptance]\label{prop:local}
Define the narrowed box $\tilde X\subseteq X^{(j+1)}$ coordinatewise: for
$i\in I_Z\cap P$ let $\tilde X_i$ be the integer hull of
$Z_i\cap X^{(j+1)}_i$; for $i\notin P$ and $G_i=\{a,a'\}$ let
$\tilde X_i=\{a'\}$ if $m_i(a)>U\ge m_i(a')$, $\tilde X_i=\{a\}$ if
$m_i(a')>U\ge m_i(a)$, and $\tilde X_i=X_i$ otherwise; and let
$\tilde X_i=X^{(j+1)}_i$ for $i\in I_C\cap P$. Let $\tilde x\in\tilde X$
with $F(\tilde x)\le U$, $\zeta=\nabla F(\tilde x)$,
$J_+=\{i:\tilde X_i\text{ has positive length}\}$, and
$\tilde X_i=[\alpha_i,\alpha_i']$ (or its integer points) for $i\in J_+$.
Suppose every $i\in J_+$ falls under one of the cases
\[
\begin{array}{lll}
\text{(i)}&\tilde x_i=\alpha_i,\ \zeta_i\ge0:&\mu_i=\zeta_i/(\alpha_i'-\alpha_i);\\
\text{(ii)}&\tilde x_i=\alpha_i',\ \zeta_i\le0:&\mu_i=-\zeta_i/(\alpha_i'-\alpha_i);\\
\text{(iii)}&\alpha_i<\tilde x_i<\alpha_i',\ \zeta_i=0:&\mu_i=0;\\
\text{(iv)}&i\in I_Z\text{ otherwise}:&\mu_i=-\abs{\zeta_i},
\end{array}
\]
and that $H_{J_+J_+}+2\diag(\mu_{J_+})\succeq0$. Then $\tilde x$ is a global
minimizer.
\end{proposition}

\begin{proof}
Let $x\in X$. If $F(x)>U$, then $F(x)>F(\tilde x)$. Otherwise
$x\in X^{(j+1)}$. Applying Lemma~\ref{lem:node}(b) to the coordinates
$i\notin P$ whose $\tilde X_i$ is a single endpoint keeps the point in
$X^{(j+1)}$ and does not increase $F$, and gives a point $x''$ that, by
Lemma~\ref{lem:node}(a), lies in $\tilde X$. Let $d=x''-\tilde x$; then
$d_i=0$ for $i\notin J_+$ and $\alpha_i-\tilde x_i\le d_i\le\alpha_i'-\tilde x_i$
for $i\in J_+$. In cases (i) and (ii),
$\zeta_id_i=\abs{\zeta_i}\abs{d_i}\ge\mu_id_i^2$; in case (iii) both sides
vanish; in case (iv) $d_i\in\Z$, so
$\zeta_id_i\ge-\abs{\zeta_i}\abs{d_i}\ge-\abs{\zeta_i}d_i^2$. Hence
$F(x'')-F(\tilde x)=\zeta^\T d+\frac12d^\T Hd\ge\frac12d_{J_+}^\T\bigl(H_{J_+J_+}
+2\diag(\mu_{J_+})\bigr)d_{J_+}\ge0$, and $F(x)\ge F(x'')\ge F(\tilde x)$.
\end{proof}

The test needs one exact $LDL^\T$ factorization of a
$\abs{J_+}\times\abs{J_+}$ matrix, which a checker repeats after replaying the
filtering history. Its validity needs neither uniqueness nor growth, and the
candidate $\tilde x$ is arbitrary. The corrected minimizer lies in
$\tilde X$: for $i\in I_Z\cap P$, $(y_j)_i$ is a node with
$m_i((y_j)_i)\le Q(y_j)=\beta_j\le\OPT\le U$, and for $i\notin P$ the same
inequality excludes the endpoint whose min-marginal exceeds $U$. The
following rule computes a candidate from $\tilde X$ and $y_j$.

\begin{definition}[Face candidate]\label{def:facecand}
Put $\tilde x_i=(y_j)_i$ for $i\in I_Z$ and for $i\notin J_+$. For
$i\in I_C\cap J_+$, put $\tilde x_i=\ell_i$ if $\ell_i\in\tilde X_i$, else
$\tilde x_i=u_i$ if $u_i\in\tilde X_i$, and otherwise call $i$ free. If the
set $J_f$ of free coordinates has $H_{J_fJ_f}$ nonsingular, the \emph{face
candidate} is the point $\tilde x$ whose free part solves
$\nabla_{J_f}F(\tilde x)=0$; otherwise stage $j$ has no face candidate.
\end{definition}

The face candidate costs one exact solve of a $\abs{J_f}\times\abs{J_f}$ linear
system. It is tested by Proposition~\ref{prop:local}, whose hypotheses
include $\tilde x\in\tilde X$ and $F(\tilde x)\le U$.

\begin{corollary}[Acceptance of the face candidate]\label{cor:local}
Let $F$ have point growth with constant $g$ at $x^*$. Let $J_0=J_0(x^*)$ and
$A=J_\partial\cap I_C^+$, the continuous coordinates with $L_i>0$ at which
$x^*$ is at a bound, and assume strict complementarity,
$\partial_iF(x^*)\ne0$ for $i\in A$. Put
$\lambda_A=\min_{i\in A}\abs{\partial_iF(x^*)}$ and
$\Gamma=\norm{H_{AA}}_2+\norm{H_{J_0A}}_2^2/(2g)$,
and let $\delta_X$ be the smallest distance $\abs{x^*_i-e}$ over coordinates
$i\in I_C\cup(I_Z\setminus P)$ and endpoints $e\in\{\ell_i,u_i\}$ with
$e\ne x^*_i$. Put $h^*=\omega^*/\sqrt{n\kappa}$, where $\omega^*$ is the
smallest of the numbers $\frac12$ (if $I_Z\cap P\ne\emptyset$), $\delta_X/6$
(if $I_C\cup(I_Z\setminus P)\ne\emptyset$) and $\lambda_A/(5\Gamma)$ (if
$A\ne\emptyset$); at least one of the first two conditions holds, because
$I_Z\cap P$ and $I_C\cup(I_Z\setminus P)$ cover $[n]$. Consider a trial
of CT with a common mesh (Lemma~\ref{lem:commonmesh}) with
$8\kappa\theta^2\le1$. At every stage $j$ with $h_j\le h^*$ that the trial
reaches and ends by filtering, the face candidate exists, equals $x^*$, and
passes the test of Proposition~\ref{prop:local}. Since $\delta_X\ge1/R$,
$\lambda_A\ge1/(\Delta R)$ and $\log_2\Gamma\le\poly(I)+\log_2\kappa$, the
least $j$ with $h_j\le h^*$ is at most $\poly(I)+O(\log\kappa)$.
\end{corollary}

The proof is in Appendix~\ref{app:localized}. It also shows that at these
stages $\tilde X_i=\{x^*_i\}$ for every integer coordinate and every
coordinate with $L_i=0$, and that $J_f=J_0$. With several minimizers the test
can fail, for example when two optimal vertices differ in a coordinate
$i\notin P$.

In experiment S1 (Section~\ref{sec:comp-localized}) the test accepted the
face candidate within nine stages on 29 of 30 random mixed-integer instances;
on the five of them on which EX used the most stages, one CT run needed 40 to
72 stages to reach the threshold of Proposition~\ref{prop:accept} with the
constants of~\eqref{eq:exact-constants}. The idea of certifying a neighborhood of a
candidate by local information is that of exclusion boxes in interval
branch-and-bound \cite{Neumaier2004}; here it is carried out exactly on the
box produced by the filtering history.


\subsection{Explicit polynomial factors}\label{sec:polynomial}

Now let each factor be an explicit list of rational monomials of degree at
most $d\ge2$, with $d$ fixed (or encoded in unary).

\begin{lemma}[Interval curvature certificate]\label{lem:intcurv}
For each $i$ write $\partial_{ii}F=\sum_mc_mx^{\alpha_m}$ and let
$[\underline\mu_m,\overline\mu_m]$ be the range of the monomial $x^{\alpha_m}$
on $\bar X$. Put $\hat L_i=\sum_m\max\{c_m\overline\mu_m,
c_m\underline\mu_m\}$ and $L_i=\max\{\hat L_i,0\}$. Then the $L_i$ are
upper coordinate curvatures, and they are computed in time polynomial in $I$.
\end{lemma}

\begin{proof}
The monomial $x^{\alpha}$ is a product of functions $x_k^{\alpha_k}$ of
distinct coordinates, so its range on the box $\bar X$ is the product of the
intervals $\{t^{\alpha_k}:t\in[\ell_k,u_k]\}$, each of which is computed
exactly from $\ell_k^{\alpha_k}$, $u_k^{\alpha_k}$ and, for even $\alpha_k$
with $\ell_k<0<u_k$, the value $0$. Hence
$c_mx^{\alpha_m}\le\max\{c_m\overline\mu_m,c_m\underline\mu_m\}$
on $\bar X$, and $\partial_{ii}F\le\hat L_i\le L_i$. For every $x\in\bar X$,
$t\mapsto F(x+te_i)-\frac{L_i}2t^2$ has second derivative
$\partial_{ii}F-L_i\le0$ on its interval, so it is concave. The endpoints are
powers of degree at most $d$ of input rationals, so their bit length is
$O(dI)$, and there are at most as many terms as input monomials.
\end{proof}

The bound can overestimate when monomials cancel; the condition number in
the results below is computed with the accepted $L$. Any sharper bound with a
certificate checkable in polynomial time may be used instead.

\begin{proposition}[Lattice stopping rule]\label{prop:lattice}
Suppose every coordinate is integer, and let $\Gamma_F$ be the product of
the denominators of the coefficients of the objective (after substituting
fixed coordinates). If $\hat x\in X$, $\beta\le\OPT$ and
$F(\hat x)-\beta<1/\Gamma_F$, then $\hat x\in\mathcal S$. It suffices that
$F(\hat x)-\beta\le2^{-q}$ with $q=1+\lceil\log_2\Gamma_F\rceil=O(I)$.
\end{proposition}

\begin{proof}
At integer points every monomial value lies in $\Gamma_F^{-1}\Z$, so
$F(X)\subseteq\Gamma_F^{-1}\Z$. If $F(\hat x)>\OPT$ then
$F(\hat x)-\OPT\ge1/\Gamma_F>F(\hat x)-\beta$, contradicting $\beta\le\OPT$.
\end{proof}

\begin{corollary}[Polynomial factors]\label{cor:poly}
For fixed $d$, under point growth~\eqref{eq:growth} with the accepted
curvature bounds, CT$(2^{-q})$ returns a certified $2^{-q}$-approximation in
$f_d(p,\kappa)(I+q+1)^{C}$ bit operations, where $f_d$ is a computable
function and $C$ an absolute constant. If all coordinates are integer,
CT$(2^{-(1+\lceil\log_2\Gamma_F\rceil)})$ returns an exact minimizer, certified
by Proposition~\ref{prop:lattice}, in $f_d(p,\kappa)(I+1)^C$ bit operations.
\end{corollary}

\begin{proof}
The interpolation, filtering, contraction, localization and grid-size
arguments of Sections~\ref{sec:grids} and~\ref{sec:growth} use only
Definition~\ref{def:curvature} and growth, not the quadratic form. In the bit
analysis of Theorem~\ref{thm:approx} (Appendix~\ref{app:growth}) only the
exponents $e_i$ and factor evaluation change. The $L_i$ of
Lemma~\ref{lem:intcurv} have $O(dI)$ bits and lie in
$(\Gamma_F\Gamma_X^{d-2})^{-1}\Z$, because the coefficients of
$\partial_{ii}F$ are integer multiples of coefficients of $F$ and the
endpoints of the monomial ranges are products of at most $d-2$ box endpoints.
So $\abs{e_i},\abs E=O(dI)$, and every node has a denominator dividing
$\Gamma_X2^\alpha$ with $\alpha=J+O(dI)+\mu K_\mu$. A monomial of degree at
most $d$ at such nodes has a value with denominator dividing
$\Gamma_F(\Gamma_X2^\alpha)^d$, and since $d\ge2$, so has $8$ times a
correction $L_iw^2/8$. Hence factor values, corrections, table entries and
messages have the common denominator $8\Gamma_F(\Gamma_X2^\alpha)^d$, and
over it they are integers of $O(d(I+\alpha))$ bits. For fixed $d$ these bit lengths exceed those of
Appendix~\ref{app:growth} by a constant factor, and the evaluation of an
explicit factor costs a polynomial in its length and these bit lengths. The
second statement follows from Proposition~\ref{prop:lattice} with $q=O(I)$.
\end{proof}

\begin{example}[Limits of the polynomial extension]\label{ex:polylimits}
(a) $F(x)=x^3-6x$ on $[1,2]$ has $\partial_{xx}F=6x\le12=\hat L$ and
$F(x)-\OPT=(x-\sqrt2)^2(x+2\sqrt2)\ge3(x-\sqrt2)^2$, so $g=3$ is a growth
constant and $\kappa\le4$; but its minimizer $\sqrt2$ and value $-4\sqrt2$
are irrational, so no rational exact output exists for continuous polynomial
coordinates. (b) $F(x)=x^4$ on $[-1,1]$ has a unique minimizer and no
$g>0$, so Lemma~\ref{lem:unique-growth} fails for polynomials.
(c) For $k\ge1$, $F(x,y)=x^2-\frac12x^{2^k}+y^2$ on $[0,1]^2$ has
$\partial_{xx}F\le2$, $\partial_{yy}F=2$ and growth constant $\frac12$, since
$x^{2^k}\le x^2$; its input has $O(k)$ bits if exponents are written in
binary, but $F(\frac12,0)$ has denominator $2^{2^k+1}$. Binary-encoded
exponents therefore need a separate evaluation analysis.
\end{example}

For continuous polynomial objectives an exact answer must therefore be
implicit. Appendix~\ref{app:boundary} gives one form: under point growth and
strict complementarity at the active bounds, sign tests on the certified
retained box identify the active face, and a uniform convexity test on the
remaining coordinates yields a strongly convex polynomial subproblem whose
unique minimizer is the global one. The cost carries a term
$B_\lambda=\lceil\log_2\max\{2,C_2/\lambda_A\}\rceil$ in the smallest inward
derivative $\lambda_A$ at the active bounds; Propositions~\ref{prop:margin}
and~\ref{prop:weakcompl} show that this term is intrinsic to certificates of
this form and that weak complementarity defeats them.
